% 电子版的紧凑排版；保持正文、摘要字号、正文行距及 25 mm 页边距。
\usepackage{placeins}
\raggedbottom
\setlength{\textfloatsep}{7pt plus 2pt minus 2pt}
\setlength{\floatsep}{7pt plus 2pt minus 2pt}
\setlength{\intextsep}{6pt plus 2pt minus 2pt}
\captionsetup{skip=4pt}

% 引导词 \hei 的定义在 common_setup.tex，供电子版与纸质版共用。

% 国赛层级：一级四号黑体 > 二级小四黑体 > 正文小四宋体 > 图号五号黑体、题文五号宋体。
\DeclareCaptionFont{fivehaohei}{\zihao{5}\heiti}
\DeclareCaptionFont{fivehaosong}{\zihao{5}\songti}
\captionsetup[figure]{labelfont=fivehaohei,textfont=fivehaosong}
\captionsetup[table]{labelfont=fivehaohei,textfont=fivehaosong}
\captionsetup[subtable]{labelfont=fivehaohei,textfont=fivehaosong}
\captionsetup[subfigure]{labelfont=fivehaohei,textfont=fivehaosong}
\AtBeginEnvironment{table}{\zihao{5}}
\AtBeginEnvironment{longtable}{\zihao{5}}

\renewcommand{\abstractname}{摘\hspace{1em}要}
\renewcommand{\thesubsection}{\arabic{section}.\arabic{subsection}}
\renewcommand{\thesubsubsection}{\thesubsection.\arabic{subsubsection}}

\setcounter{topnumber}{4}
\setcounter{bottomnumber}{3}
\setcounter{totalnumber}{6}
\renewcommand{\floatpagefraction}{0.65}

\makeatletter
% 摘要结束后会 \newpage\null。\null 在下页顶占住一行，节标题前的弹性空白就不会被丢掉，
% 「一、问题重述」会被空出一大截。去掉 \null，标题贴页眉下 25 mm 正文区顶。
\patchcmd{\endabstract}
  {\endquotation\newpage\null}
  {\endquotation\newpage}
  {}{\PackageWarning{layout_setup}{Could not remove the blank box after the abstract.}}

% 标题保持模板原字体（三号加粗），单行居中。

% 浮动专页从顶端排起，避免把几张图表之间的空隙拉长。
\setlength{\@fptop}{0pt}
\setlength{\@fpsep}{8pt}
\setlength{\@fpbot}{0pt plus 1fil}

% 收紧标题空白；一级四号黑体，二级/三级小四黑体，保证小节标题重于表题表头。
\renewcommand\section{\FloatBarrier\@startsection{section}{1}{\z@}%
  {-2.0ex plus -.5ex minus -.2ex}{1.0ex plus .2ex}%
  {\centering\normalfont\zihao{4}\heiti\bfseries}}
\renewcommand\subsection{\@startsection{subsection}{2}{\z@}%
  {-1.7ex plus -.4ex minus -.2ex}{.7ex plus .2ex}%
  {\normalfont\zihao{-4}\heiti\bfseries}}
\renewcommand\subsubsection{\@startsection{subsubsection}{3}{\z@}%
  {-1.4ex plus -.3ex minus -.2ex}{.5ex plus .2ex}%
  {\normalfont\zihao{-4}\heiti\bfseries}}
\renewcommand\paragraph{\@startsection{paragraph}{4}{\z@}%
  {1.4ex plus .3ex minus .2ex}{-1em}%
  {\normalfont\zihao{-4}\heiti\bfseries}}

% 保留表内字体，减少模板默认 1.38 倍的额外行高。
\renewenvironment{tabular}
  {\bgroup\renewcommand{\arraystretch}{1.12}\mcm@oldtabular}
  {\mcm@endoldtabular\egroup}

% 文献采用模板的小号字（约 11 pt），条目间不额外插入空行。
\patchcmd{\thebibliography}{\list}{\small\list}
  {}{\PackageError{layout_setup}{Bibliography font patch failed}{Check the bibliography definition.}}
\patchcmd{\thebibliography}{\sloppy}
  {\setlength{\itemsep}{0pt}\setlength{\parsep}{0pt}\sloppy}
  {}{\PackageError{layout_setup}{Bibliography spacing patch failed}{Check the bibliography definition.}}

% 公式与正文间留出清晰间隔，避免每条短公式额外占据两行。
\g@addto@macro\normalsize{%
  \abovedisplayskip=7pt plus 2pt minus 2pt
  \belowdisplayskip=7pt plus 2pt minus 2pt
  \abovedisplayshortskip=3pt plus 1pt
  \belowdisplayshortskip=5pt plus 1pt minus 1pt}
\makeatother
